\section{Pre-specified hypotheses}\label{app:design}
The following were recorded in \texttt{docs/02\_research\_design.md} before the evaluation runs.
\begin{itemize}
\item \textbf{H1.} \LX{} achieves higher HR@1 and a lower normalised rank than tag-level attributions aggregated post hoc.
\item \textbf{H4.} \LX{} beats exact Shapley over data-driven clusters (CLX) and random partitions (RPX) with the same number of groups.
\item \textbf{H2.} \TEA{} increases HR@$k$ at a fixed delivery time and increases temporal stability.
\item \textbf{H3.} \LX+\TEA{} has deterministic cost, lower latency than sampling explainers at equal or higher accuracy, and the highest DCL and in-time rate.
\item \textbf{H5a.} Localization degrades less under sensor noise with \LX+\TEA.
\item \textbf{H5b.} A decoy captures top-1 less often with \LX+\TEA.
\item \textbf{H6.} Random window splits inflate detection metrics relative to chronological, attack-disjoint splits.
\end{itemize}
A hypothesis is reported as supported only if the Holm-corrected two-sided $p<0.05$ and the effect has the predicted sign.

\section{Validation protocols in the reviewed corpus}\label{app:corpus}
\begin{table}[h]
\centering\scriptsize
\caption{Validation protocols stated in reviewed XAI-IDS studies that use time-series CPS/IIoT data (as stated in the papers; ``--'' if not stated in the sections we read). Random proportional splits or $k$-fold CV without a stated chronological or gapped protocol allow overlapping or adjacent windows on both sides of the split.}\label{tab:corpus}
\begin{tabular}{@{}lll@{}}
\toprule
Study & Data & Stated protocol\\
\midrule
\citet{ha2022ifac} & gas pipeline & 80/20 split\\
\citet{yazdinejad2025ppfl} & SWaT, gas pipeline, UNSW-NB15 & 70/30 split\\
\citet{birihanu2024correlation} & SWaT, HAI, Modbus & 80/20 split\\
\citet{xu2024alrl} & SWaT, WADI & five-fold CV\\
\citet{rasheed2026xaicps} & SWaT, WADI, BATADAL, IIoT & 5-fold CV, 70/15/15\\
\citet{panesar2026mmdcritic} & HIL microgrid scalograms & 70/30 split\\
\citet{shoukat2025trustmyids} & N-BaIoT, Edge-IIoTset & 70/30 split\\
\citet{attique2024iiot} & CICIDS2017, X-IIoTID & 70/30 split\\
\citet{saiyed2026eshield} & WUSTL-IIoT-2021, CICIoT2023 & stratified 5-fold CV\\
\citet{franco2026shats} & SWaT & padded segment split (leakage-aware)\\
\citet{fung2024attributions} & SWaT, WADI, TEP & normal-only training, attack test sets\\
\bottomrule
\end{tabular}
\end{table}

\section{HAI 21.03 loop map}\label{app:loops}
The ten groups used for HAI 21.03 are listed below. Tags constant in training were removed, leaving 57 tags in nine non-empty groups; P1-AUX holds discrete pump-state tags, all constant. The full mapping, with engineering-judgement flags, is in \texttt{configs/hai/hai2103\_loops.yaml}, and the ground truth is in \texttt{configs/hai/hai2103\_attacks.csv}.
\begin{itemize}
\item P1-PC: pressure;
\item P1-LC: return-tank level;
\item P1-FC: discharge flow;
\item P1-TC: heat-exchanger temperature;
\item P1-AUX: pumps;
\item P2-SC: turbine speed;
\item P2-TC: vibration and trip limits;
\item P3-LC: upper-tank level;
\item P4-ST: steam-turbine model;
\item P4-HT: hydro-turbine model.
\end{itemize}

\section{Additional results}\label{app:extra}
\begin{table}[htbp]
\centering\scriptsize
\caption{Loop HR@1 per detector family (``at alarm, single-shot / 60\,s after the alarm with \TEA''; mean over attacks of the per-attack mean over 3 seeds; $n$: attacks detected by at least one seed).}\label{tab:perdet}
\resizebox{\linewidth}{!}{%
\begin{tabular}{@{}lllccccccc@{}}
\toprule
Data & Detector & $n$ & RE & IG & KernelSHAP-1024 & Fung ens. & PI-CAX & \LX-bg & \LX\\
\midrule
HAI & linear & 46 & 0.70 / 0.66 & 0.85 / 0.78 & 0.82 / 0.72 & 0.87 / 0.82 & 0.93 / 0.93 & 0.85 / 0.80 & 0.85 / 0.80\\
 & mlp & 45 & 0.75 / 0.64 & 0.71 / 0.70 & 0.70 / 0.70 & 0.73 / 0.72 & 0.93 / 0.93 & 0.71 / 0.68 & 0.75 / 0.68\\
 & cnn & 45 & 0.73 / 0.60 & 0.67 / 0.69 & 0.70 / 0.68 & 0.70 / 0.70 & 0.94 / 0.93 & 0.70 / 0.65 & 0.75 / 0.67\\
 & gru & 45 & 0.66 / 0.53 & 0.76 / 0.73 & 0.74 / 0.65 & 0.72 / 0.64 & 0.92 / 0.91 & 0.70 / 0.70 & 0.81 / 0.63\\
\midrule
PCL-Sim & linear & 270 & 0.89 / 0.80 & 0.89 / 0.87 & 0.87 / 0.76 & 0.88 / 0.82 & 0.89 / 0.88 & 0.89 / 0.86 & 0.89 / 0.84\\
 & mlp & 226 & 0.87 / 0.85 & 0.87 / 0.86 & 0.86 / 0.85 & 0.90 / 0.87 & 0.94 / 0.88 & 0.85 / 0.86 & 0.85 / 0.83\\
 & cnn & 253 & 0.80 / 0.76 & 0.79 / 0.76 & 0.78 / 0.78 & 0.81 / 0.76 & 0.85 / 0.84 & 0.83 / 0.79 & 0.79 / 0.76\\
 & gru & 265 & 0.87 / 0.81 & 0.88 / 0.87 & 0.86 / 0.81 & 0.88 / 0.84 & 0.89 / 0.87 & 0.89 / 0.88 & 0.89 / 0.86\\
\bottomrule
\end{tabular}}
\end{table}

\begin{table}[htbp]
\centering\scriptsize
\caption{Tag-level localization (HR@1 / HR@3, single-shot at the given time after the alarm). H-\LX{} refines the top loop of \LX{} into tags; the other methods rank all tags directly.}\label{tab:tags}
\resizebox{\linewidth}{!}{%
\begin{tabular}{@{}llcccc@{}}
\toprule
Data & Method & +0\,s & +30\,s & +60\,s & +120\,s\\
\midrule
HAI & RE & 0.23 / 0.40 & 0.13 / 0.27 & 0.13 / 0.27 & 0.02 / 0.11\\
 & SM & 0.10 / 0.34 & 0.13 / 0.32 & 0.13 / 0.26 & 0.10 / 0.20\\
 & IG & 0.28 / 0.41 & 0.13 / 0.27 & 0.15 / 0.29 & 0.07 / 0.15\\
 & KS-1024 & 0.31 / 0.41 & 0.15 / 0.29 & 0.16 / 0.27 & 0.08 / 0.17\\
 & LIME-1024 & 0.31 / 0.41 & 0.15 / 0.28 & 0.17 / 0.29 & 0.06 / 0.16\\
 & Fung-ens & 0.36 / 0.48 & 0.23 / 0.37 & 0.22 / 0.36 & 0.10 / 0.19\\
 & H-\LX & 0.31 / 0.49 & 0.12 / 0.34 & 0.15 / 0.31 & 0.06 / 0.18\\
\midrule
PCL-Sim & RE & 0.38 / 0.76 & 0.23 / 0.37 & 0.26 / 0.43 & 0.26 / 0.42\\
 & SM & 0.20 / 0.53 & 0.18 / 0.37 & 0.17 / 0.38 & 0.16 / 0.37\\
 & IG & 0.44 / 0.74 & 0.21 / 0.42 & 0.25 / 0.47 & 0.22 / 0.42\\
 & KS-1024 & 0.44 / 0.73 & 0.14 / 0.31 & 0.19 / 0.37 & 0.18 / 0.33\\
 & LIME-1024 & 0.44 / 0.73 & 0.15 / 0.32 & 0.21 / 0.40 & 0.20 / 0.36\\
 & Fung-ens & 0.28 / 0.75 & 0.14 / 0.36 & 0.19 / 0.41 & 0.18 / 0.38\\
 & H-\LX & 0.45 / 0.76 & 0.17 / 0.32 & 0.21 / 0.39 & 0.19 / 0.36\\
\bottomrule
\end{tabular}}
\end{table}

\begin{table}[htbp]
\centering\scriptsize
\caption{PCL-Sim loop HR@1 by attack primitive (at alarm, single-shot / +60\,s with \TEA; $n$: detected (attack, detector) units).}\label{tab:prim}
\resizebox{\linewidth}{!}{%
\begin{tabular}{@{}llcccc@{}}
\toprule
Primitive & $n$ & RE & IG & \LX-bg & \LX\\
\midrule
CV & 135 & 0.87 / 0.82 & 0.86 / 0.84 & 0.87 / 0.82 & 0.88 / 0.83\\
CV\_PV & 160 & 0.96 / 0.91 & 0.95 / 0.88 & 0.95 / 0.88 & 0.95 / 0.87\\
CV\_ST & 159 & 0.94 / 0.96 & 0.93 / 0.96 & 0.94 / 0.98 & 0.93 / 0.96\\
PV\_BIAS & 148 & 0.93 / 0.85 & 0.94 / 0.95 & 0.93 / 0.94 & 0.93 / 0.93\\
PV\_FREEZE & 115 & 0.56 / 0.54 & 0.51 / 0.49 & 0.56 / 0.54 & 0.50 / 0.46\\
SP & 137 & 0.86 / 0.82 & 0.90 / 0.89 & 0.90 / 0.87 & 0.90 / 0.86\\
SP\_PV & 160 & 0.79 / 0.66 & 0.81 / 0.80 & 0.86 / 0.81 & 0.81 / 0.76\\
\bottomrule
\end{tabular}}
\end{table}

